\subsection{特征值变分刻画的应用}

在本节中，我们考虑 C 上的希尔伯特空间。

命题 7. --- 设 u 与 v 是 E 的正紧自同态。

\begin{enumerate}
\def\labelenumi{\alph{enumi})}
\item
  有 \(|\lambda_n(u) - \lambda_n(v)|\leqslant \| u - v\|\) （对所有 \(n\in \mathbf{I}\) ）；
\item
  若 \(u \leqslant v\) ，则 \(\lambda_n(u) \leqslant \lambda_n(v)\) （对所有 \(n \in I\) ）。
\end{enumerate}

设 \(n \in I\) ，设 F 是 E 的 n 维向量子空间。对所有 \(x \in F\) ，我们有

\[
| \langle x | v (x) \rangle - \langle x | u (x) \rangle | \leqslant \| u - v \| \| x \| ^ {2},
\]

因此有不等式

\[
\mathrm{R} _ {\mathrm{F}} (v) - \| u - v \| \leqslant \mathrm{R} _ {\mathrm{F}} (u) \leqslant \mathrm{R} _ {\mathrm{F}} (v) + \| u - v \|.
\]

断言 \(a\) ) 由此按 IV, p. 154 命题 6 得出。

若 \(u \leqslant v\) ，我们有 \(\langle x | u(x) \rangle \leqslant \langle x | v(x) \rangle\) （对所有 \(x \in E\) ）。因此对每个 \(n \in I\) 与每个子空间 \(F\) （维数为 \(n\) 的），我们有 \(R_F(u) \leqslant R_F(v)\) ，从而 \(\lambda_n(u) \leqslant \lambda_n(v)\) （同处）。

命题 8. --- 设 u 是 E 的正紧自同态。设 H 是 E 的一个闭子空间，\(i_{\mathrm{H}}: \mathrm{H} \to \mathrm{E}\) 是典范单射。用 \(u_{\mathrm{H}}\) 表示 H 的自同态 \(i_{\mathrm{H}}^{*}ui_{\mathrm{H}}\) 。它是紧且正的。

\begin{enumerate}
\def\labelenumi{\alph{enumi})}
\item
  有 \(\mathrm{I_H}\subset \mathrm{I_E}\) 与 \(\lambda_n(u_{\mathrm{H}})\leqslant \lambda_n(u)\) （对所有 \(n\in \mathrm{I_H}\) ）；
\item
  若 H 在 E 中有有限余维数 \(k \in N\) ，则我们有 \(I_{H} + k \subset I_{E}\) 与 \(\lambda_{n+k}(u) \leqslant \lambda_{n}(u_{\mathrm{H}})\) （对所有 \(n \in I_{H}\) ）。
\end{enumerate}

自同态 \(u_{\mathrm{H}}\) 是紧的（III, p. 5, 命题 3）。它是正的，因为 \(\langle x \mid u_{\mathrm{H}}(x) \rangle = \langle i_{\mathrm{H}}(x) \mid u(i_{\mathrm{H}}(x)) \rangle \geqslant 0\) （对所有 \(x \in \mathrm{H}\) ）。

设 \(n \in I_{H} \subset I_{E}\) 。设 F 是 H 的一个维数为 \(n + 1\) 的适应于 \(u_{H}\) 的子空间。于是我们有 \(\lambda_{n}(u_{\mathrm{H}}) = r_{\mathrm{F}}(u_{\mathrm{H}})\) （IV, p. 153, 命题 5, a)），又由于 \(r_{\mathrm{F}}(u_{\mathrm{H}}) = r_{\mathrm{F}}(u) \leqslant \lambda_{n}(u)\) （IV, p. 154, 命题 6），我们得到断言 a)。

假设 H 在 E 中有余维数 \(k \in N\) 且 \(n \in I_{H}\) 。设 F 是 H 的一个维数为 n 的适应于 \(u_{H}\) 的子空间。它在 H 中的正交等于 \(H \cap F^{\circ}\) ，且它是 E 中子空间 \(F + H^{\circ}\) （维数为 \(n + k\) ）的正交。因此 \(n + k \in I_{E}\) ，且（IV, p. 153, 命题 5, b)）

\[
\lambda_{n}(u_{\mathrm{H}}) = \sup_{\substack{x\in \mathrm{H}\cap \mathrm{F}^{\circ}\\ x\neq 0}}\frac{\langle x\mid u_{\mathrm{H}}(x)\rangle}{\|x\|^{2}} = \mathrm{R}_{\mathrm{F} + \mathrm{H}^{\circ}}(u)
\]

从而 \(\lambda_{n}(u_{\mathrm{H}})\leqslant \lambda_{n + k}(u)\) （IV, p. 154, 命题 6）。

定义 3. --- 设 F 是希尔伯特空间，u 是从 E 到 F 中的一个紧线性映射。

对每个整数 \(n \in \mathrm{I}_{\mathrm{E}}\) ，置 \(\alpha_{n}(u) = \lambda_{n}(u^{*} \circ u)\) 。设 J 是由 \(n \in \mathrm{I}_{\mathrm{E}}\) （使 \(\alpha_{n}(u) > 0\) ）组成的集。族 \((\alpha_{n}(u))_{n \in \mathrm{J}}\) 称为 \(u\) 的按重数重复的奇异值序列。

我们称序列 \((\alpha_{n}(u))_{n\in\mathrm{I}_{\mathrm{E}}}\) 为 u 的增广奇异值序列。

序列 \((\alpha_{n}(u))_{n\in\mathrm{I}_{\mathrm{E}}}\) 是完全定义的，因为 E 的自同态 \(u^{*}\circ u\) 是紧的（III, p. 5, 命题 3）；由于 \(u^{*}\circ u\) 是正的，它是一个正实数的递减族。

命题 9. --- 设 F 是希尔伯特空间，u 是从 E 到 F 中的一个紧线性映射。

\begin{enumerate}
\def\labelenumi{\alph{enumi})}
\tightlist
\item
  对 \(n \in \mathrm{I}_{\mathrm{E}}\) ，我们有
\end{enumerate}

\[
\alpha_ {n} (u) = \sup _ {\mathrm{F} \in \mathscr {F} _ {n + 1}} \inf _ {x \in \mathrm{F} - \{0 \}} \frac {\| u (x) \|}{\| x \|} = \inf _ {\mathrm{F} \in \mathscr {F} _ {n}} \sup _ {x \in \mathrm{F} ^ {\circ} - \{0 \}} \frac {\| u (x) \|}{\| x \|}.
\]

\begin{enumerate}
\def\labelenumi{\alph{enumi})}
\setcounter{enumi}{1}
\tightlist
\item
  设 J 是由 \(n \in I_{E}\) （使 \(\alpha_{n}(u) \neq 0\) ）组成的集。存在 E 中的标准正交族 \((e_{n})_{n \in \mathrm{J}}\) 与 F 中的标准正交族 \((f_{n})_{n \in \mathrm{J}}\) ，使得对每个 \(x \in E\) ，我们有
\end{enumerate}

\[
u (x) = \sum_ {n \in J} \alpha_ {n} (u) \langle e _ {n} | x \rangle f _ {n}.
\]

由于 \(\langle x \mid u^{*}u(x)\rangle = \|u(x)\|^{2}\) （对所有 \(x \in E\) ），诸定义与 IV, p. 154, 命题 6, a) 蕴含第一个断言中的等式。设 \((e_{n})_{n\in\mathrm{I}_{\mathrm{E}}}\) 是 E 的一个标准正交族，满足把 IV, p. 152 命题 3 应用于 E 的紧正自同态 \(u^{*} \circ u\) 所得的结论。置 \(f_{n} = \alpha_{n}^{-1}u(e_{n})\) （对 \(n \in J\) ）。按 IV, p. 150 推论 2 的证明中的推理，我们得到断言 c)。

推论. --- 设 F 是希尔伯特空间，u 是从 E 到 F 中的一个紧线性映射。设 \(w \in \mathcal{L}(\mathrm{E})\) 与 \(v \in \mathcal{L}(\mathrm{F})\) 。对每个 \(n \in I_{E}\) ，我们有 \(\alpha_{n}(w \circ u \circ v) \leqslant \|v\|\) \(\|w\|\) \(\alpha_{n}(u)\) 。

这是前一命题断言 \(a\) ) 的推论。

注记. --- 1) 序列 \((\alpha_{n}(u))_{n\in\mathrm{I}_{\mathrm{E}}}\) 是递减的，故集 J 或等于 \(\mathrm{I}_{\mathrm{E}}\) ，或等于线段 \(\{0,\ldots,m\}\) （\(\mathrm{I}_{\mathrm{E}}\) 中的）（其中 \(m\in\mathbf{N}\) ）。后一情形成立当且仅当 \(u\) 有限秩。

\begin{enumerate}
\def\labelenumi{\arabic{enumi})}
\setcounter{enumi}{1}
\item
  由于 \(\| u(x)\| = \| |u|(x)\|\) （对所有 \(x\in \mathrm{E}\) ；I, p. 139, 命题 10），我们有 \(\alpha_{n}(u) = \alpha_{n}(|u|)\) （对所有 \(n\in \mathrm{I}_{\mathrm{E}}\) ）。
\item
  若 \(u\) 是正的，则 \(\alpha_{n}(u)=\lambda_{n}(u)\) （对所有 \(n\in\mathrm{I}_{\mathrm{E}}\) ；事实上，按 IV, p. 153 命题 4，我们这时有 \(\alpha_{n}(u)^{2}=\lambda_{n}(u^{*}u)=\lambda_{n}(u^{2})=\lambda_{n}(u)^{2}\) ）。
\item
  有可能有 \(\alpha_{n}(v\circ u\circ w)=0\) ，即使 \(\alpha_{n}(u)\) 非零；在这种情形，\(\alpha_{n}(v\circ u\circ w)\) 不是 \(v\circ u\circ w\) 的奇异值。
\end{enumerate}
% semantic-fragments: legacy-input-seam
\relax{}
